\chapter{Análisis de la Industria y la Competencia}\label{ch:competitive-analysis}

% Alcance: las cinco fuerzas, la cadena de valor, estructura del sector y rentabilidad (HHI).
% Enlace cruzado: la ventaja duradera que sobrevive a estas fuerzas es el tema de
%   \cref{ch:advantage-and-disruption}; el poder de fijación de precios trazado aquí tiene su
%   raíz en \cref{ch:pricing} (índice de Lerner) y \cref{ch:customer-economics} (LTV).

La rentabilidad de una industria no es aleatoria. Las condiciones estructurales que rodean a
una empresa --- el apalancamiento de los proveedores, el apalancamiento de los compradores,
los sustitutos, los nuevos entrantes y la rivalidad interna --- establecen un techo a los
rendimientos antes de que la estrategia siquiera comience. La cadena de valor pregunta,
entonces, dónde dentro de esa industria la empresa captura lo que queda.

\section{Las Cinco Fuerzas de Porter}\label{sec:five-forces}
% Complejidad: 4 → figura five-forces

Antes de asignar capital o entrar a un mercado, el estratega necesita saber qué características
estructurales limitarán los rendimientos y cuáles los apoyarán. Las Cinco Fuerzas responden esa
pregunta antes de que se firme un solo cliente. El mecanismo es directo: un mercado de alto margen
atrae entrantes y compradores poderosos hasta que los rendimientos se comprimen al costo de capital.
Las Cinco Fuerzas nombran los cinco canales por los cuales ocurre esa compresión y preguntan qué
tan intensa es cada una en este momento --- no si la industria es atractiva en abstracto, sino si
\emph{esta empresa} puede sostener una posición en ella.

El \emph{framework} es cualitativo; su trabajo analítico es una secuencia diagnóstica aplicada a
cada fuerza por separado \parencite{porter1980}:
\begin{description}
  \item[Rivalidad entre competidores existentes] Concentración (véase \cref{eq:hhi}),
        transparencia de precios, diferenciación del producto, costos de salida e incrementos
        de capacidad. Los costos fijos elevados y los productos de tipo \emph{commodity} hacen
        intensa la rivalidad.
  \item[Amenaza de nuevos entrantes] Escala mínima eficiente, tecnología propia,
        efectos de red, lealtad de marca, costos de cambio y barreras regulatorias.
        Los bajos requerimientos de capital y la tecnología madura reducen esta barrera de entrada.
  \item[Amenaza de sustitutos] Productos de categorías adyacentes que cumplen el mismo
        trabajo-por-hacer; evaluados mediante la elasticidad cruzada de precios (\cref{eq:elasticity}).
        Los sustitutos establecen un techo de precio suave.
  \item[Poder de negociación de compradores] Concentración de compradores, costos de cambio,
        amenaza de integración hacia atrás, participación del gasto del comprador y sensibilidad
        al precio.
  \item[Poder de negociación de proveedores] Concentración de proveedores, unicidad del insumo,
        amenaza de integración hacia adelante y costo de cambio de proveedor.
\end{description}

\begin{figure}[htbp]
  \centering
  \includegraphics[width=0.86\linewidth]{five-forces}
  \caption{Las Cinco Fuerzas. La rivalidad competitiva ocupa el centro; cuatro fuerzas ejercen
    presión hacia adentro sobre la rentabilidad del sector. La intensidad de cada flecha
    representa la magnitud de la negociación o de la amenaza.}
  \label{fig:five-forces}
\end{figure}

\begin{admonition}{Advertencia}
Las Cinco Fuerzas son una fotografía estática a nivel de industria. Omiten tres aspectos cada
vez más relevantes en los mercados tecnológicos. Primero, los \emph{complementos de plataforma}:
la rentabilidad de un mercado digital depende de la inversión de los complementadores, que no son
ni proveedores ni compradores en el sentido tradicional. Segundo, la \emph{competencia dinámica}:
el modelo no tiene reloj; una fuerza que hoy parece débil (por ejemplo, la entrada) puede ser
irrelevante en un mercado que se reestructura con rapidez. Tercero, la visión basada en recursos
(\cref{ch:advantage-and-disruption}) sostiene que la heterogeneidad de capacidades a nivel de
empresa --- no sólo la estructura del sector --- explica las diferencias de rendimiento dentro
de una industria. Utilizar las Cinco Fuerzas para elegir la arena; utilizar VRIN para preguntar
si se puede ganar en ella.
\end{admonition}

\begin{example}
Aplicación al mercado B2B SaaS de referencia. \emph{Decisión: ¿conviene invertir en
diferenciación o competir en precio?}

\begin{itemize}
  \item \textbf{Poder de compradores --- alto.} Los compradores son pymes con bajos costos de
        cambio (las exportaciones de datos son estándar), múltiples alternativas y un gasto por
        asiento reducido en relación con sus presupuestos. Esto comprime el poder de fijación de
        precios hacia el piso del índice de Lerner (\cref{eq:lerner}).
  \item \textbf{Amenaza de entrada --- alta.} La infraestructura sin código y las APIs abiertas
        permiten ensamblar una solución puntual creíble por bastante menos de \money{500000}.
        Los efectos de red son débiles al momento del lanzamiento.
  \item \textbf{Rivalidad --- intensa.} El mercado está fragmentado (HHI $\approx$\num{2150},
        véase \cref{eq:hhi}); la paridad de funcionalidades es habitual; los sitios de comparación
        de precios hacen visible el cambio de proveedor. A \money{50}/mes el producto está plenamente
        en la zona de peligro de \emph{commodity}.
  \item \textbf{Sustitutos --- moderados.} Las hojas de cálculo y herramientas SaaS más simples
        sirven al mismo trabajo; la amenaza es real pero manejable si las integraciones generan
        \emph{stickiness}.
  \item \textbf{Poder de proveedores --- bajo.} La infraestructura en la nube (AWS, GCP) es
        competitiva; ningún insumo propietario está cautivo en un solo proveedor.
\end{itemize}

El alto poder de compradores más la alta amenaza de entrada en un mercado fragmentado impone una
respuesta de diferenciación --- singularidad de producto que eleve los costos de cambio --- en
lugar de una carrera de liderazgo en costos que la empresa no puede ganar a su escala actual.
\end{example}

\begin{admonition}{Advertencia}
Tratar las cinco fuerzas como si tuvieran el mismo peso y promediarlas en una única
«puntuación de atractivo». Un mercado puede ser \emph{estructuralmente poco atractivo} en
cuatro dimensiones y aun así ser muy rentable para una empresa que haya neutralizado la quinta
(por ejemplo, una barrera de patentes que elimina la amenaza de entrada). Evaluar cada fuerza
por separado; luego preguntar cuál de ellas aborda la estrategia.
\end{admonition}

\textcite{dranove7theconomics} amplían las Cinco Fuerzas con un análisis de la rivalidad
basado en la teoría de juegos: cuando los competidores interactúan de forma repetida, las
estrategias de ojo-por-ojo pueden sostener precios supracompetitivos sin necesidad de colusión.
El modelo también se inscribe dentro del paradigma más amplio «estructura--conducta--desempeño»
(SCP, por sus siglas en inglés), cuya validez empírica es objeto de debate (véase
\cref{ch:advantage-and-disruption} para la crítica basada en recursos).

El diagnóstico de las Cinco Fuerzas establece el contexto para cada decisión de precios en
\cref{ch:pricing} y cada elección de segmentación en
\cref{ch:segmentation-targeting-positioning}: entrar sólo donde al menos una fuerza sea
genuinamente manejable. \textcite{porter1980} y \textcite{dranove7theconomics} desarrollan el
sustento empírico de la economía industrial.

\section{La Cadena de Valor}\label{sec:value-chain}
% Complejidad: 3

¿Qué actividades dentro de la empresa generan el costo o la diferenciación que importa al
cliente, y cuáles son gastos generales que deben minimizarse o externalizarse? Una empresa no
es sólo un producto; es una secuencia de actividades. Tanto el costo como la diferenciación
se originan en eslabones específicos de esa cadena. Saber qué eslabón importa es la condición
previa para asignar la inversión correctamente --- y para entender dónde un competidor con una
estructura de cadena diferente puede socavar o superar a la empresa.

\textcite{porter1985} organiza las actividades de la empresa en dos niveles:
\begin{description}
  \item[Actividades primarias] Las que tocan directamente el producto o servicio:
    logística de entrada $\to$ operaciones $\to$ logística de salida $\to$
    \emph{marketing} y ventas $\to$ servicio postventa.
  \item[Actividades de apoyo] Las que habilitan las primarias: infraestructura de la empresa,
    gestión de recursos humanos, desarrollo tecnológico y adquisiciones.
\end{description}
Para una empresa B2B SaaS, la cadena primaria se comprime a cuatro eslabones:
desarrollo de producto $\to$ ventas y \emph{marketing} $\to$ éxito del cliente $\to$
infraestructura/operaciones. Las actividades de apoyo (finanzas, legal, RRHH) se aplican sin
cambios. El «margen» es el excedente agregado que queda después del costo combinado de todos
los eslabones.

La cadena de valor \emph{no} es una cadena de suministro: las cadenas de suministro rastrean
insumos y productos físicos entre empresas; la cadena de valor mapea actividades que agregan
valor \emph{dentro} de una sola empresa. Confundirlas conduce a decisiones de externalización
basadas únicamente en el costo, ignorando la diferenciación y el conocimiento tácito
incorporados en la actividad. El modelo también asume que las actividades son separables, lo
que sobreestima la separabilidad en productos digitales altamente integrados.

\begin{example}
\emph{Decisión: ¿dónde gasta la empresa SaaS y dónde está construyendo ventaja?}

Mezcla observada de gastos operativos (aproximación, coherente con un margen bruto
$\approx$\qty{70}{\percent}):
\begin{itemize}
  \item Ventas y \emph{marketing}: $\approx$\qty{50}{\percent} del ingreso
        (el eslabón dominante --- coherente con un presupuesto publicitario de \money{120000}/mes
        más un equipo de ventas).
  \item I+D / producto: $\approx$\qty{30}{\percent} del ingreso.
  \item Éxito del cliente / infraestructura: $\approx$\qty{15}{\percent} del ingreso.
  \item G\&A: $\approx$\qty{5}{\percent}.
\end{itemize}
Las ventas y el \emph{marketing} dominan porque la empresa está en la etapa pre-\emph{moat}:
debe comprar clientes que aún no puede retener de forma orgánica. Una empresa post-\emph{moat}
(\cref{ch:growth}) desplaza el gasto hacia producto y éxito del cliente a medida que los costos
de cambio reducen el costo marginal de retención. La lente de la cadena de valor hace visible
esta trayectoria como un desplazamiento en qué eslabón lleva la inversión.
\end{example}

\begin{admonition}{Advertencia}
Equiparar la cadena de valor con un desglose de costos por actividad y luego «optimizar» cada
eslabón de forma independiente. El valor de la cadena proviene de los \emph{vínculos}: una
actividad de servicio superior aumenta el valor de la actividad de producto al reducir el
\emph{churn} (véase \cref{eq:retention} en \cref{ch:customer-economics}). Recortar el área
de éxito del cliente para reducir el costo por eslabón destruye el valor del vínculo ---
y eleva el \emph{churn}.
\end{admonition}

La cadena de valor es anterior al modelo de negocio de plataforma, donde las actividades
primarias se invierten: la plataforma produce una infraestructura de intermediación y el valor
se crea en las interacciones entre los participantes de terceros, no en las operaciones propias
de la empresa. Las implicaciones para la estructura de costos y el análisis competitivo se
exploran en \cref{ch:business-models}.

La concentración de costos de la cadena de valor en ventas y \emph{marketing} refleja la
estructura competitiva que las Cinco Fuerzas diagnosticaron antes: cuando la amenaza de entrada
y el poder de compradores son ambos altos, la demanda debe adquirirse de forma costosa. El
camino hacia la mejora del margen pasa por construir el \emph{moat} primero (\cref{ch:growth})
y luego reconfigurar la cadena. \textcite{porter1985} es la fuente primaria.

\section{Estructura del Sector y Rentabilidad}\label{sec:industry-structure}
% Complejidad: 3

¿Qué tan concentrado es el mercado en el que competimos, y qué implica eso sobre los niveles
de precio sostenibles y la intensidad competitiva? Un mercado con cuatro rivales aproximadamente
iguales se comporta de manera diferente a uno con una empresa dominante y una franja periférica.
La concentración de mercado no es destino --- un mercado concentrado puede ser ferozmente
competitivo (véanse las aerolíneas) --- pero es el primer hecho estructural que hay que
establecer. El Índice de Herfindahl--Hirschman (HHI) traduce las cuotas de mercado en un
único número que tanto reguladores como estrategas utilizan.

Sea \(i\) una empresa con cuota \(s_i \in [0,1]\) (expresada como fracción). El HHI es:
\begin{equation}\label{eq:hhi}
  \mathrm{HHI} \;=\; \sum_{i} s_i^2.
\end{equation}
Expresado en cuotas en puntos porcentuales (es decir, \(s_i\) en \([0,100]\)), el HHI varía
desde cerca de 0 (atomístico) hasta \num{10000} (monopolio puro). Las guías antimonopolio
de EE.UU. \parencite{dranove7theconomics} señalan mercados por encima de \num{2500} como
altamente concentrados y fusiones que eleven el HHI en más de \num{200} en un mercado por
encima de \num{1500}.

La \cref{eq:hhi} asume que las cuotas de mercado aproximan la intensidad competitiva, pero dos
empresas con 50\%\ cada una pueden coludirse tácitamente o competir de forma brutal. El índice
también es sensible a cómo se define el mercado: una definición amplia deflacta el HHI y puede
ocultar la dominancia local. El índice no dice nada sobre las barreras de entrada ni las
asimetrías de costos; es un punto de partida, no una conclusión.

\begin{example}
\emph{Decisión: ¿está nuestro mercado lo suficientemente concentrado como para apoyar el
liderazgo en precios, o debemos competir en valor?}

Supóngase un mercado con cuatro actores con cuotas
\(s = (0.30,\,0.25,\,0.20,\,0.15)\) y una franja competitiva que suma \(0.10\):
\[
  \begin{aligned}
    \mathrm{HHI} &= 0.30^2 + 0.25^2 + 0.20^2 + 0.15^2 + 0.10^2 \\
                 &= 0.090 + 0.063 + 0.040 + 0.023 + 0.010 \\
                 &= 0.226 \;\approx\; \num{2260}\text{ (pp-cuadrado)}.
  \end{aligned}
\]
Con \num{2260}, el mercado está \emph{moderadamente concentrado} --- por debajo del umbral de
alta concentración de \num{2500}. El liderazgo en precios es débil; la empresa líder no puede
subir precios sin arriesgar pérdida de cuota. Una \emph{startup} que entre con
\(s \approx 0.02\) aporta \(0.0004\) al HHI --- insignificante para el agregado, pero no
para su propio estado de resultados. La entrada es viable precisamente porque no se ha
alcanzado el precipicio de concentración: ninguna empresa tiene la escala para responder de
forma aplastante.
\end{example}

\begin{admonition}{Advertencia}
Usar el HHI en solitario para concluir que un mercado es «competitivo». Un mercado
fragmentado de 50 empresas con \qty{2}{\percent} de cuota cada una (HHI = \num{200}) puede
seguir siendo \emph{estructuralmente poco atractivo} si los costos de entrada son nulos y los
sustitutos son abundantes --- las Cinco Fuerzas dicen lo que el HHI no puede decir.
\end{admonition}

El índice de Lerner (\cref{eq:lerner}) conecta el HHI con los márgenes precio-costo a través
de la elasticidad precio promedio de la industria: en un oligopolio de Cournot simétrico,
\(\mathcal{L} = \mathrm{HHI}/|\varepsilon|\). Esto vincula directamente la estructura de
mercado con el poder de fijación de precios y es el puente entre la organización industrial
y \cref{ch:pricing}. \parencite{mankiw2020,dranove7theconomics}.

La concentración de mercado cuantifica la arena; la cadena de valor revela dónde dentro de
ella la empresa crea y captura valor. Ambas alimentan el análisis de ventaja en
\cref{ch:advantage-and-disruption}. \textcite{mankiw2020} y \textcite{dranove7theconomics}
fundamentan la literatura empírica de organización industrial detrás de \cref{eq:hhi}.
